\section{Reflexão e o que quebra no nativo}
\label{sec:reflexao-nativo}

No modo JVM, reflexão custa mais que uma chamada direta; no nativo,
reflexão não declarada é um bug. A análise estática só inclui o que é
alcançável, e o que não foi alcançado não entra no binário. A página
\citetitle{doc:graalvm-metadata} explica onde está a fronteira: o Native
Image \enquote{will compute reachability at build time given that all
the call arguments are constant}, e prefere essa forma de fornecer
metadados. É o caso de \texttt{porReflexao} na seção~\ref{sec:gizmo},
que usa a constante \texttt{"privado"}. O problema real são nomes que
só existem em runtime, como os de um container que varre o classpath ao
subir. A mesma página diz ainda que, para ler ou escrever campos por
reflexão, os nomes dos campos precisam estar nos metadados; sem isso,
\texttt{Field.get} e \texttt{Field.set} lançam
\texttt{MissingReflectionRegistrationError}. O
\citetitle{doc:graalvm-compat} lista os recursos dinâmicos que exigem
metadados no build.

\subsection*{O que costuma quebrar}

O guia \citetitle{doc:quarkus-native-tips} traz os casos típicos:
\begin{description}
\item[Serialização de classes não registradas] Com Jackson, o sintoma
  citado é \texttt{No serializer found for class \ldots{} and no
  properties discovered}; com JSON-B, o JSON pode sair completamente
  vazio. O guia explica a regra geral: elementos usados via reflexão
  \enquote{are not part of the call tree so they are dead code
  eliminated}, a menos que sejam registrados.
\item[Recursos não declarados] Arquivos do classpath que devem estar no
  executável precisam ser configurados explicitamente; o caminho mais
  fácil, segundo o guia, é \texttt{quarkus.native.resources.includes}.
\item[Inicialização em build time] O guia afirma: \enquote{By default,
  Quarkus initializes all classes at build time}. Junte isso à
  seção~\ref{sec:augmentation}: o \texttt{<clinit>}
  (seção~\ref{sec:class-loading}) roda na máquina de build, e os campos
  \texttt{static} vão para o binário. Um \texttt{static final} que lê
  uma variável de ambiente guarda o valor que ela tinha na máquina de
  build. A saída é
  \texttt{--initialize-at-run-time=<pacote ou classe>}.
\item[Membros \texttt{private} em beans] Não quebram, mas custam: o guia
  \citetitle{doc:quarkus-cdi-reference} explica que acessá-los exige
  reflexão, e cada registro de reflexão aumenta o executável. Campos
  injetados \texttt{private} já são reescritos como
  \emph{package-private} por padrão (seção~\ref{sec:gizmo}).
\item[Busca programática de beans] Um bean usado só via
  \texttt{CDI.current()} pode ter sido removido no build
  (seção~\ref{sec:jandex}).
\end{description}

\subsection*{Como registrar}

\begin{itemize}
\item \texttt{@RegisterForReflection} na classe (ou numa classe de
  configuração, com \texttt{targets}, quando a classe é de uma
  biblioteca). O guia avisa que, por padrão, a anotação registra também
  classes aninhadas.
\item Um arquivo \texttt{reflect-config.json} em
  \texttt{src/main/resources/META-INF/native-image/<group-id>/<artifact-id>/},
  lido automaticamente pelo build nativo, para casos que a anotação não
  cobre; o guia avisa que, nesse caminho, manter o arquivo atualizado entre
  versões do Mandrel/GraalVM é responsabilidade sua.
\item Dentro de extensões, build items como \texttt{ReflectiveClassBuildItem}
  (guia \citetitle{doc:quarkus-writing-extensions}). É por isso que
  bibliotecas com extensão Quarkus \enquote{simplesmente funcionam} no
  nativo: alguém já escreveu esse registro.
\end{itemize}

\subsection*{A decisão}

\begin{description}
\item[Quando aceitar reflexão] Modo JVM, código de borda pouco
  executado, biblioteca madura sem alternativa. Meça antes de otimizar.
\item[Quando eliminar reflexão] Alvo nativo, caminho quente, ou quando
  existe uma extensão Quarkus para a biblioteca. Prefira injeção pelo
  construtor e campos \emph{package-private}, e escolha bibliotecas com
  extensão.
\item[O que medir] Rode a suíte de \texttt{@QuarkusIntegrationTest}
  contra o executável nativo no CI (seção~\ref{sec:testes-camadas}):
  é o teste que mostra reflexão faltando antes da produção (o
  \emph{tracing agent} do GraalVM é outro caminho, não coberto aqui). No modo JVM, o JFR e um
  profiler (seção~\ref{sec:ferramentas}) mostram se
  \texttt{Method.invoke} ou \texttt{Field.get} aparece no caminho quente.
\end{description}
